\documentclass[11pt]{article}
\usepackage[margin=1in]{geometry}
\usepackage{enumitem}
\usepackage{hyperref}
\usepackage{longtable}
\usepackage{booktabs}

\title{Nudge v0.1 System Specification}
\author{Tahaseen Shaik}
\date{Foundation revision: August 29, 2026}

\begin{document}
\maketitle

\section{Status and authority}
This document is the authoritative design source for Nudge v0.1. The target
completion date is September 16, 2026. The present revision specifies the
repository foundation and interface contracts; it does not assert validated
research hypotheses or experimental results. Items marked TODO are unresolved
and shall not be silently assigned implementation values.

\section{Purpose and scope}
Nudge is a ROS~2 Jazzy/C++ corrective shared-autonomy system for a USAR-inspired
differential-drive UGV operating in partially known environments. A human
operator provides sparse corrections through an Xbox controller. Nudge tests
whether input communicates disagreement at four semantic levels:
\texttt{MOTION}, \texttt{PLAN}, \texttt{GOAL}, and \texttt{AUTHORITY}.
\texttt{UNKNOWN} represents an uninterpreted correction.

Version 0.1 compares an explicit labeling condition with an implicit heuristic
condition. Both conditions shall publish the same interpretation contract and
use the same downstream assistance, autonomy, arbitration, safety, and logging
architecture. Learned classification, reinforcement learning, camera
perception, natural-language interaction, and physical-robot dependencies are
out of scope for this revision.

\section{Research and application context}
The primary research question is how an autonomous robot should interpret and
respond to sparse physical corrective input. Secondary questions concern
semantic-level inference, legibility, the roles of autonomy and interpretation
confidence, interaction history, unequal environmental knowledge, operator
location, and explicit-versus-implicit accuracy, effort, and performance.

The following are provisional hypotheses, not established claims: short
corrections may indicate motion adjustments; repeated corrections after route
changes may indicate planning or objective mismatch; strong command
disagreement may justify more human authority; repeated unresolved disagreement
shall force human-dominant control regardless of confidence; confidence sources
should remain separate; and joint uncertainty may require conservative
exploration rather than control transfer alone.

The application is USAR-inspired reconnaissance in a partially mapped damaged
structure. Route selection is represented conceptually by
\[
J(\tau)=w_gJ_{goal}+w_rJ_{risk}+w_tJ_{traversability}+w_eJ_{efficiency}
       +w_cJ_{context}+w_nJ_{nudge}.
\]
Terms and weights may be deterministic and configuration-driven in v0.1, but
their definitions and values are TODO. A nudge updates preferences from the
current state. A global replan begins at the current robot pose and retains the
current assumed goal.

Nudge is an engineering and HRI prototype grounded in prior shared-control,
goal-inference, corrective-interaction, policy-blending, and confidence-based
arbitration work. Its design directions shall not be described as validated
novelty before evaluation and literature review support that claim.

\section{Architecture constraints}
The core stack shall be simulator-agnostic and shall not call Gazebo-specific
APIs. Gazebo integration belongs in \texttt{nudge\_sim}; a future Isaac Sim
backend and physical UGV shall be replaceable backends. Standard ROS~2
interfaces shall be used when their semantics suffice. The safety layer shall
remain independent of both human and autonomous command generation. Operator
context shall be represented and logged from v0.1 even while behavior remains
context-independent.

\section{Enumerated contracts}
The following numeric encodings are stable interface values.

\begin{longtable}{lll}
\toprule
Concept & Name & Value \\
\midrule
\endhead
Correction type & UNKNOWN, MOTION, PLAN, GOAL, AUTHORITY & 0, 1, 2, 3, 4 \\
Interpretation source & EXPLICIT, HEURISTIC & 0, 1 \\
Operator context & ADJACENT, ONBOARD, REMOTE & 0, 1, 2 \\
Authority state & AUTONOMOUS, SHARED, HUMAN\_DOMINANT & 0, 1, 2 \\
Lifecycle state & INIT, READY, RUNNING, PAUSED, COMPLETE, ESTOP & 0, 1, 2, 3, 4, 5 \\
\bottomrule
\end{longtable}

\section{ROS interface fields}
All state and event messages carry \texttt{std\_msgs/Header} unless they are
embedded records. Identifiers are strings to avoid selecting a UUID policy in
v0.1.

\subsection{CorrectionEvent}
Fields: header; \texttt{event\_id}; \texttt{active};
\texttt{geometry\_msgs/Twist human\_command}; normalized
\texttt{magnitude}; \texttt{duration\_ms}; optional explicit semantic label
encoded as \texttt{explicit\_label} plus \texttt{has\_explicit\_label}; and
\texttt{operator\_context}. Magnitude normalization is TODO and shall be
configured rather than assumed.

\subsection{CorrectionInterpretation}
Fields: header; \texttt{correction\_type}; \texttt{confidence};
\texttt{interpretation\_source}; \texttt{matched\_heuristic\_rule\_id}; and
the complete \texttt{source\_event}. Because the wire field is unsigned,
zero denotes no heuristic match (including explicit interpretations), while
H01--H99 map to their numeric suffix. Confidence calibration and UNKNOWN
handling are TODO.

\subsection{AssistanceState}
Fields: header; \texttt{authority\_state}; \texttt{lifecycle\_state};
\texttt{human\_weight}; \texttt{autonomy\_weight};
\texttt{safety\_override\_active}; \texttt{transition\_reason}; and
\texttt{unresolved\_disagreement\_count}. Blend-weight policy is TODO; the
fields make the applied arbitration observable.

\subsection{ConfidenceState}
Fields: header; \texttt{autonomy\_confidence};
\texttt{correction\_confidence}; and per-value validity flags. The confidence
definition, estimator, calibration, and update rate are TODO.

\subsection{OperatorState}
Fields: header; \texttt{operator\_context}; \texttt{debug\_mode}; cumulative
\texttt{correction\_count}; \texttt{unresolved\_disagreement\_count}; and
\texttt{cumulative\_intervention\_ms}. Debug mode is strictly an observability
flag and shall not transition the authority FSM.

\subsection{RobotGoalState}
Fields: header; \texttt{has\_assumed\_goal};
\texttt{geometry\_msgs/PoseStamped assumed\_goal}; \texttt{goal\_id}; and
monotonic \texttt{revision}. Goal selection behavior is out of scope.

\section{Controller contract}
Axis and button IDs shall be configurable. The left stick provides corrective
motion. LT holds or increases human authority; RT yields or increases autonomous
authority. A accepts/confirms, B rejects the current plan, X pauses, and Y
requests a global replan from the current pose to the current assumed goal. The
D-pad supplies labels in Explicit Semantic Nudge mode. View/Back toggles debug
mode. Menu/Start requests emergency stop. Platform-specific axis polarity and
trigger representation are TODO configuration decisions.

The proposed explicit mapping is D-pad Up to MOTION, Right to PLAN, Down to
GOAL, and Left to AUTHORITY. It remains configurable pending usability review.

\section{Heuristic worksheet contract}
The implicit condition is deterministic and rule-based in v0.1. It may use
duration, normalized magnitude, human-autonomy directional disagreement,
decision-point context, repetition, and alignment with alternate goals. Raw
features shall remain logged when rules discretize them. IDs H01--H02 are
reserved for candidate MOTION rules, H03--H04 for PLAN, H05--H06 for GOAL,
H07--H08 for AUTHORITY, and H09 for UNKNOWN. All remain disabled until exact
conditions are reviewed.

Candidate features include
\[
m_H=\sqrt{x_H^2+y_H^2},\quad m_H\in[0,1],
\qquad
\Delta\theta=\cos^{-1}\left(
\frac{u_H\cdot u_A}{\lVert u_H\rVert\lVert u_A\rVert}\right)
\]
for nonzero command vectors. Zero-vector handling is TODO. The source worksheet
lists candidate values only: deadzone 0.08--0.10, start 0.18--0.20, release
0.08--0.10, release timeout 150~ms, short maximum 400~ms, long minimum
1000~ms, strong magnitude 0.75, disagreement 90--120 degrees, repeat window
2500--4000~ms, repeat limit 3, branch radius 1.5~m, alternate-goal alignment
30--45 degrees, high correction confidence 0.75, autonomy return 0.60--0.75,
and recovery 1000--2000~ms. YAML records them as comments while runtime values
remain \texttt{null}.

\section{State machines}
Authority states are \texttt{AUTONOMOUS}, \texttt{SHARED}, and
\texttt{HUMAN\_DOMINANT}. Adjacent transitions are allowed in both directions;
direct \texttt{HUMAN\_DOMINANT} to \texttt{AUTONOMOUS} is forbidden. Repeated
unresolved disagreement forces \texttt{HUMAN\_DOMINANT} independently of
correction confidence. Eventual return to autonomy is confidence-based and must
pass through \texttt{SHARED}. Threshold and timing values are TODO YAML values.

Lifecycle states are \texttt{INIT}, \texttt{READY}, \texttt{RUNNING},
\texttt{PAUSED}, \texttt{COMPLETE}, and \texttt{ESTOP}. Valid nominal
transitions are INIT--READY, READY--RUNNING, RUNNING--PAUSED,
PAUSED--RUNNING, and RUNNING--COMPLETE. ESTOP is reachable from any state.
Reset/recovery from ESTOP and restart behavior after COMPLETE are deliberately
TODO safety decisions, so the foundation FSM rejects them.

\section{Control, confidence, and assistance}
Planned conditions are MANUAL, AUTO, FIXED\_BLEND, NUDGE\_EXPLICIT, and
NUDGE\_IMPLICIT. Manual applies $u=u_H$, autonomous applies $u=u_A$, and fixed
blend applies $u=\alpha u_H+(1-\alpha)u_A$ for configured $\alpha$. A Nudge
policy may select a discrete action from autonomy confidence $C_A$, correction
confidence $C_C$, spatial environmental confidence $C_E$, observable operator
history $M_H$, and state $s$; it is not restricted to one blend value. Repeated
disagreement always selects human-dominant authority, and the safety layer
overrides unsafe human or autonomous output.

Candidate autonomy-confidence inputs include route-cost separation,
feasibility, obstacle clearance, environmental knowledge, and goal ambiguity.
Candidate correction-confidence inputs include duration, magnitude, direction
consistency, decision timing, repetition, command disagreement, and
alternate-goal alignment. Definitions, normalization, calibration, decay, and
persistence are TODO. Environmental confidence remains spatial,
$C_E(x,y)\in[0,1]$; future work may separate robot and human knowledge as
$C_{E,R}(x,y)$ and $C_{E,H}(x,y)$.

\section{Operator, environment, and observability}
The operator model records observable behavior rather than subjective mental
state: intervention count/frequency, mean duration/magnitude, consistency,
response latency, autonomy rejection, task success, and preferred assistance.
In v0.1 these values are primarily logged and operator context does not change
assistance. Candidate epistemic states are known safe, known unsafe, robot
unknown/human known, robot known/human unknown, and jointly unknown; responses
to them are not yet implemented.

Debug output should expose lifecycle and authority, assumed goal, global and
local plans, human/autonomous/executed commands, separate confidence values,
predicted and explicit labels, objective weights, active nudge cost, recent
disagreement, and matched rule ID. Debug mode remains authority-neutral.

\section{ROS package and topic boundaries}
The packages are \texttt{nudge\_interfaces}, \texttt{nudge\_ugv\_description},
\texttt{nudge\_input}, \texttt{nudge\_corrections},
\texttt{nudge\_autonomy}, \texttt{nudge\_arbitration},
\texttt{nudge\_operator}, \texttt{nudge\_environment},
\texttt{nudge\_logging}, \texttt{nudge\_sim}, and
\texttt{nudge\_bringup}. Candidate topics are \texttt{/joy},
\texttt{/human/cmd}, \texttt{/autonomy/cmd},
\texttt{/autonomy/confidence}, \texttt{/environment/confidence},
\texttt{/correction/event}, \texttt{/correction/interpretation},
\texttt{/operator/state}, \texttt{/robot/assumed\_goal},
\texttt{/assistance/state}, and the final safe \texttt{/cmd\_vel}.

\section{Configuration}
Controller mapping and input filtering, event detection, classifier thresholds,
confidence thresholds, FSM thresholds, and logging options shall live in YAML.
Unfinalized values shall be represented by YAML \texttt{null} with a TODO
comment. Production nodes shall validate required parameters before entering
READY rather than substituting hidden defaults.

\section{Robot description}
The nominal \texttt{nudge\_ugv} is a parameterized differential-drive robot
with an approximately 0.45 by 0.32 by 0.16~m body, 11~kg nominal body mass,
0.08~m wheel radius, 0.30~m wheel separation, rear caster, 0.8~m/s maximum
linear speed, 1.5~rad/s maximum angular speed, 2D LiDAR, and optional IMU. These
are engineering scaffold values, not experimental findings. The TF chain is
map--odom--base\_footprint--base\_link with wheel, caster, lidar, and IMU child
links. Dynamic map/odom transforms are supplied by localization/navigation, not
the URDF.

\section{Logging contract}
Each run directory is named \texttt{<timestamp>\_<scenario\_id>} and contains
\texttt{manifest.json}, \texttt{telemetry.csv}, \texttt{events.jsonl}, and
\texttt{rosbag2/}. Control conditions are MANUAL, AUTO, FIXED\_BLEND,
NUDGE\_EXPLICIT, and NUDGE\_IMPLICIT. Required manifest fields are run ID,
timestamp, git commit, scenario ID, control condition, operator context, robot
assumed goal, ground-truth goal, intended correction type, objective weights,
confidence thresholds, controller configuration, and map ID.

Event types are CORRECTION\_START, CORRECTION\_END,
CORRECTION\_CLASSIFIED, PLAN\_COST\_UPDATED, GOAL\_ASSUMPTION\_CHANGED,
GLOBAL\_REPLAN, AUTHORITY\_TRANSITION, REPEATED\_DISAGREEMENT,
UNKNOWN\_REGION\_ENTERED, SAFETY\_OVERRIDE, and TASK\_COMPLETE. Classified
correction records additionally preserve matched heuristic rule ID, optional
explicit label, inferred label, correction confidence, operator context,
robot assumed goal, and objective weights before and after.

Telemetry includes time, robot pose/velocity, human/autonomous/executed
commands, all three confidence sources, authority state, assumed goal, active
plan ID, and the human-input-active flag. The format should remain compatible
with future interaction flight-recorder work.

\section{Safety, scenarios, and evaluation}
The independent safety filter has final authority, including over human input.
Candidate constraints include velocity, acceleration, clearance, restricted
regions, emergency stop, and stale-command timeout; exact values and reset
behavior require review.

Initial scenarios isolate local motion correction, route rejection, goal
correction, authority takeover, low environmental confidence, and optionally
joint uncertainty. Controlled scenario metadata supplies semantic ground truth
and operator instruction; the implicit classifier shall never receive that
label.

Candidate metrics include completion time, path length/cost, human input,
intervention/route-update/replan counts, takeover frequency, success,
classification error/accuracy, unsafe states, autonomy utilization, and time in
each authority state. The explicit/implicit comparison should at least examine
accuracy, completion time, effort, intervention count, unnecessary authority
transitions, and incorrect persistent route/goal changes. Human-subject metrics
are deferred pending a reviewed study design.

\section{Implementation priorities and scope risks}
The planned sequence is robot/simulator/Nav2 baseline, controller,
manual/auto/fixed-blend baselines, FSMs, event extraction, explicit semantics,
worksheet finalization, heuristic interpretation, current-state objective
updates, confidence, logging/debugging, scenarios/ground truth, and evaluation.
This foundation implements only repository structure, interfaces,
configuration, minimal FSMs, robot description, logging schema, and tests.

Scope risks include excessive simulation fidelity, learned inference before
deterministic baselines, treating heuristics as findings, conflating observable
behavior with subjective confidence, category proliferation, uncontrolled
environment-uncertainty scope, premature hardware integration, and Gazebo
coupling. Later versions may add learned interpretation, persistent adaptation,
knowledge asymmetry, richer legibility, Isaac Sim, physical deployment,
operator-context studies, human-subject evaluation, and flight-recorder or
explanation integration.

\section{Open design decisions}
TODO items include all heuristic and confidence thresholds, controller device
IDs/polarities, magnitude normalization, objective definitions and weights,
confidence estimators and calibration, UNKNOWN behavior, repeated-disagreement
resolution semantics, safety reset from ESTOP, COMPLETE restart semantics,
scenario design, evaluation metrics, and experimental protocol. These require
explicit decisions in later revisions.

\end{document}
